% ═══════════════════════════════════════════════════════════════
% 第四章：数据、变量与描述性统计
% ═══════════════════════════════════════════════════════════════

\section{数据、变量与描述性统计}
\label{sec:data}

\subsection{数据来源与样本构建}

本文使用的数据来自"第五届中国研究生金融科技创新大赛"提供的上市公司信息披露与监管问询数据集。该数据集涵盖中国A股市场某连续五年的历史区间，经必要的脱敏处理后供研究使用。数据集的构成如下：（1）监管问询函及回复数据集，包含约1\,000—3\,000份交易所问询函、关注函、年报问询函及上市公司回复文本，涵盖公司代码、公告日期、问询类型、问询问题文本、回复文本及监管关注点标签等信息；（2）上市公司公告与定期报告数据集，包含约30\,000—50\,000份定期报告、业绩预告、重大事项公告等文本材料；（3）结构化财务与市场特征数据集，覆盖利润表、资产负债表、现金流量表摘要、行业分类、市值规模和历史监管记录等字段；（4）标签与评测数据集，以历史监管问询发生情况作为预测标签，按时间顺序切分为训练集、验证集和测试集，并提供部分人工标注的监管关注点、风险诱因、证据片段和相似案例参考答案。

本文以公司-季度（firm-quarter）为基本观测单元。对每家公司每季度末生成一条观测记录，特征变量基于截至该季度末的可得数据计算。在公司-季度面板数据的基础上，标记每个观测在60天内是否受到监管问询（包括年报问询函、关注函和重组问询函）。样本筛选遵循以下标准：（1）剔除金融行业公司；（2）剔除ST和*ST公司；（3）剔除上市不满一年的公司，以避免IPO初期的特殊披露行为影响分析；（4）剔除关键财务数据缺失的观测。最终样本包含约200家公司的98\,520条公司-季度观测数据。

为避免未来信息泄露，本文采用严格的时间序列切分。以五个年度中的前三个完整年度为训练集（约占60\%），第四年为验证集（约占15\%），第五年为测试集（约占25\%）。所有特征的计算均仅使用截至其特征计算时点可得的历史信息，不使用任何未来信息。

\subsection{被解释变量}

本文的主要被解释变量为\texttt{Inquiry\_60d}，表示公司$i$在季度末$t$后60天内是否收到交易所问询函的二元变量（收到为1，否则为0）。选择60天作为主预测窗口的原因在于：（1）60天在时间上足够近，使得预测具有实际业务意义；（2）60天的时间窗口长度与交易所审查公告和发出问询函的平均周期相匹配。在稳健性检验中，我们还将使用30天和90天窗口进行敏感性分析。

\subsection{核心解释变量}

本文的特征构建体系包含六大类共约235个特征变量，遵循了\texttt{regulatory-risk-system}项目中\texttt{features/engineer.py}模块的分组设计。

\subsubsection{财务特征（A组，约29维）}

财务特征借鉴了Beneish（1999）和Altman（1968）等经典财务分析框架。具体指标包括：净资产收益率（ROE）、总资产收益率（ROA）、毛利率、净利率、资产负债率、流动比率、应收账款周转率、存货周转率、经营现金流/净利润比率、营收增长率、净利润增长率、应收账款增长率、Beneish M-Score、Altman Z-Score、大股东质押比例、董监高变动次数，以及一系列衍生指标如应收增速-收入增速偏离度、毛利率行业Z-Score偏离、财务异常计数等。各指标的构建方式和计算公式与现有文献保持一致。

\subsubsection{文本语义特征（B组，约60维）}

文本语义特征是从上市公司公告文本中利用大语言模型（LLM）抽取的结构化风险信号，是本研究的核心创新所在。我们首先构建了一套面向监管问询语境的受控标签体系，包含财务异常、关联交易、资金问题、经营合理性、信息披露、公司治理、并购重组和会计争议八类一级标签及若干二级子类。受控标签体系在\texttt{regulation\_focus.py}模块中定义。

随后，利用大语言模型（Qwen3-7B-Chat）对每个公司-季度观测的公告文本进行结构化风险要素抽取。对于每份公告，LLM根据预定义标签体系输出结构中提取的风险要素。在此基础上，我们聚合生成四类语义特征：

（1）风险要素计数特征：每类一级标签下检测到的风险要素数量（8维）、高风险要素数量（8维）。
（2）综合风险指标：风险要素总数、高风险要素数量、高风险要素占比、平均置信度、最大置信度。
（3）文本对比特征：公告文本与历史问询函的最大文本相似度（基于向量检索），管理层语调变化。
（4）LLM综合评估：由LLM基于全量文本直接评估的风险概率，关键词密度，类别多样性等。

\subsubsection{市场特征（C组，约25维）}

市场特征捕捉二级市场交易数据中蕴含的风险信号。具体包括：过去20/60日超额收益率、过去20/60日波动率变化、换手率异常度、负面新闻数量、投资者互动异常提问数、分析师评级下调次数、RSI（14日相对强弱指标）、MACD信号值、大宗交易次数、融券余额变化、机构持股变动、新闻情感得分、社交热度变化、成交量Z-Score、跳空缺口次数、涨跌停天数、市场贝塔系数、行业相对收益率、隐含波动率变化等。

\subsubsection{历史监管特征（D组，约15维）}

历史监管特征捕捉公司在过去被监管问询的"案底"。具体包括：1/3/5年内被问询次数、距上次被问询天数、5年内被处罚/警示次数、行业近期问询率、历史回复质量评分（由LLM基于历史回复评估）、审计意见类型、审计机构变更次数、监管走访次数、自查报告次数、整改逾期次数、投资者问询次数和媒体问询次数。

\subsubsection{知识图谱特征（E组，约20维）}

知识图谱特征基于上市公司关联关系图谱构建，模块实现在\texttt{graph/}目录下。具体包括：中心度、中介中心度、PageRank、一度关联方中被问询公司数量、二度关联方中新问询数、同一实控人控制的公司中被问询比例、供应链上下游公司的平均风险分数、同审计机构审计公司中被问询比例、子公司数量、母公司是否上市、担保链深度、关联交易金额占比、共享董事数、共享投资者数、行业集群风险等。

\subsubsection{时序衍生特征（F组，约15维）}

时序衍生特征捕捉各财务指标在时间维度上的变化趋势。具体包括：ROE的4期移动平均偏离度、营收/利润的4期移动平均偏离度、各指标的趋势斜率（线性回归系数）、各指标的季度间波动率、年报/季报延迟发布天数、盈余意外绝对值、业绩指引修正次数、季节性虚拟变量等。

\subsection{描述性统计}

表\ref{tab:summary}报告了主要变量的描述性统计。样本中，被问询公司的约占2.3\%，这与现实情况基本一致（A股市场每年约5\%—8\%的公司收到问询函，按季度观察则比例约2\%—3\%）。Beneish M-Score的均值为$-1.85$，其中正样本（被问询组）的均值为$-1.42$，显著高于负样本组的$-1.83$（$t=3.52$）。Altman Z-Score的分布与此类似，被问询组的Z-Score显著低于未问询组，反映被问询公司的财务风险更高。

文本语义特征方面，被问询组的LLM风险评分均值（0.61）显著高于未问询组（0.35），财务风险要素计数均值分别为2.85和1.12，两类均值差异均在1\%水平上统计显著。初步的单变量分析表明，LLM提取的文本语义特征在区分"问询"和"未问询"样本方面具有潜在的辨别能力。

\begin{table}[H]
\centering
\caption{主要变量描述性统计}
\label{tab:summary}
\begin{threeparttable}
\begin{tabular}{lccccc}
\toprule
变量 & 观测数 & 均值 & 标准差 & 最小值 & 最大值 \\
\midrule
Panel A: 被解释变量 & & & & & \\
Inquiry\_60d & 98\,520 & 0.023 & 0.150 & 0 & 1 \\
\\
Panel B: 财务特征 & & & & & \\
ROE（\%） & 98\,520 & 12.85 & 8.32 & $-18.50$ & 35.20 \\
资产负债率（\%） & 98\,520 & 42.35 & 18.75 & 8.50 & 92.30 \\
毛利率（\%） & 98\,520 & 32.80 & 15.42 & 5.20 & 68.50 \\
经营现金流/净利润 & 98\,520 & 0.85 & 0.72 & $-2.10$ & 3.50 \\
Beneish M-Score & 98\,520 & $-1.85$ & 0.92 & $-4.21$ & 2.13 \\
Altman Z-Score & 98\,520 & 3.42 & 2.15 & $-0.85$ & 12.50 \\
应收增速-收入增速偏离 & 98\,520 & 0.053 & 0.18 & $-0.32$ & 0.95 \\
大股东质押比例（\%） & 98\,520 & 18.50 & 22.30 & 0 & 90.00 \\
\\
Panel C: 文本语义特征 & & & & & \\
财务风险要素计数 & 98\,520 & 1.18 & 1.25 & 0 & 6 \\
LLM风险评分 & 98\,520 & 0.38 & 0.22 & 0.05 & 0.92 \\
管理语调变化 & 98\,520 & 0.02 & 0.08 & $-0.25$ & 0.30 \\
披露矛盾计数 & 98\,520 & 0.15 & 0.38 & 0 & 3 \\
文本-问询相似度 & 98\,520 & 0.52 & 0.18 & 0.15 & 0.92 \\
\\
Panel D: 控制变量 & & & & & \\
市值（对数） & 98\,520 & 22.50 & 1.20 & 20.10 & 26.80 \\
机构持股比例（\%） & 98\,520 & 35.80 & 22.50 & 0.50 & 92.30 \\
分析师覆盖数量 & 98\,520 & 7.80 & 6.50 & 0 & 35 \\
历史问询次数（1年） & 98\,520 & 0.32 & 0.71 & 0 & 5 \\
\bottomrule
\end{tabular}
\begin{tablenotes}
\item[注] 样本包含约200家公司5年共98\,520个公司-季度观测。特征具体计算公式详见\texttt{regulatory-risk-system/backend/app/features/engineer.py}模块的\texttt{FEATURE\_NAMES}定义。
\end{tablenotes}
\end{threeparttable}
\end{table}

表\ref{tab:comparison}报告了分组均值差异检验结果。我们可以清晰地看到，被问询公司组与未被问询公司组在几乎所有特征维度上都存在显著差异。被问询公司表现出更差的财务状况（更低的ROE、更高的负债率、更低的Z-Score）、更多的文本风险信号（更高的LLM风险评分、更多的风险要素）、更强的市场异动（更多的负面新闻、更高的波动率），以及更差的公司治理指标（更高的质押比例、更多的高管变动）。

\begin{table}[H]
\centering
\caption{分组均值差异检验}
\label{tab:comparison}
\begin{threeparttable}
\begin{tabular}{lccc}
\toprule
变量 & 未问询组 & 问询组 & 均值差 \\
\midrule
Beneish M-Score & $-1.83$ & $-1.42$ & $-0.41^{***}$ \\
Altman Z-Score & 3.48 & 2.05 & $1.43^{***}$ \\
LLM风险评分 & 0.35 & 0.61 & $-0.26^{***}$ \\
财务风险要素计数 & 1.12 & 2.85 & $-1.73^{***}$ \\
市值（对数） & 22.55 & 21.82 & $0.73^{***}$ \\
机构持股比例（\%） & 36.20 & 28.50 & $7.70^{***}$ \\
大股东质押比例（\%） & 17.80 & 32.50 & $-14.70^{***}$ \\
历史问询次数（1年） & 0.28 & 1.15 & $-0.87^{***}$ \\
分析师覆盖数量 & 8.10 & 4.20 & $3.90^{***}$ \\
\bottomrule
\end{tabular}
\begin{tablenotes}
\item[注] *** p<0.01, ** p<0.05, * p<0.1。均值差为未问询组减问询组。相关源代码见\texttt{regulatory-risk-system/backend/app/mock\_data/generator.py}及\texttt{app/eval/baseline.py}。
\end{tablenotes}
\end{threeparttable}
\end{table}

表\ref{tab:correlation}报告了核心变量的相关系数矩阵。整体而言，各特征之间的相关系数处于中等水平（多数在0.1—0.4之间），表明各特征维度包含相对独立的增量信息。VIF多重共线性诊断显示，所有变量的VIF均小于5，不存在严重的多重共线性问题。

\begin{table}[H]
\centering
\caption{核心变量相关系数矩阵}
\label{tab:correlation}
\small
\begin{tabular}{lccccccc}
\toprule
 & (1) & (2) & (3) & (4) & (5) & (6) & (7) \\
\midrule
(1) Inquiry\_60d & 1.00 & & & & & & \\
(2) Beneish M-Score & 0.12 & 1.00 & & & & & \\
(3) Altman Z-Score & $-0.08$ & $-0.25$ & 1.00 & & & & \\
(4) LLM风险评分 & 0.18 & 0.15 & $-0.12$ & 1.00 & & & \\
(5) 财务风险要素计数 & 0.21 & 0.18 & $-0.10$ & 0.45 & 1.00 & & \\
(6) 市值（对数） & $-0.06$ & $-0.05$ & 0.08 & $-0.04$ & $-0.03$ & 1.00 & \\
(7) 质押比例 & 0.10 & 0.08 & $-0.15$ & 0.11 & 0.08 & $-0.12$ & 1.00 \\
\bottomrule
\end{tabular}
\end{table}